% =====================================================================
% Paper 53 (standalone): spectral truncation of the flat plane --
% state-space convergence with a squared window -- and the finite
% disk-with-cone left open.
%
% REBUILT 2026-09-29 (PI-approved).  What this paper now claims:
%   * [INTERNAL THEOREM, on the named state classes; canonicity open] Flat plane R^2,
%     band-limited window, translation seminorm:  the compression /
%     lifted-state pattern of Paper 38 (after Gaudillot-Estrada--van
%     Suijlekom) with a POSITIVE squared-window kernel |psi|^2 gives
%     d_GH <= gamma_Lambda = gamma_1 Lambda^{-1/2} between the
%     Monge--Kantorovich state spaces (finite-first-moment classes).
%   * [MEASURED] The 2D Bochner--Riesz kernel is negative somewhere at
%     every tested order and its Lebesgue constant exceeds 1; this is why
%     the squared window, not Bochner--Riesz, is the lift.
%   * [MEASURED] Finite disk: the Markov--Cesaro map is NOT positivity
%     preserving (ring counterexample); its reconstruction error on the
%     open disk DECAYS on every tested observable.
% What was withdrawn in the rebuild:
%   * Bochner--Riesz positivity (false); the plane "propinquity" via a
%     Bochner--Riesz tunnel (not established); any backbone for Paper 51's
%     cigar-tip entropy (the order-0 Hankel "plane" is the flat radial
%     sector -- alpha never entered; Paper 51 cites this paper only for
%     the withdrawn material);  the finite-disk "non-decaying rate Lambda^{+0.07}" (a
%     floating-point artifact: the 0/0 Markov ratio evaluated exactly at
%     rho = R);  the boundary-collar error 0.26 at R = 12 (not
%     reproduced);  finite-disk Cesaro positivity as a general property.
% Backing: tests/test_paper53_flat_plane_squared_window.py,
%   tests/test_paper53_plane_bochner_riesz.py,
%   tests/test_paper53_disk_obstruction.py,
%   tests/test_paper53_height_constant.py.
% Siblings: Paper 38 (SU(2)), 40 (compact Lie groups).
% =====================================================================
\documentclass[11pt,a4paper]{article}

\usepackage[T1]{fontenc}
\usepackage[utf8]{inputenc}
\usepackage{amsmath, amssymb, amsthm, mathrsfs, mathtools}
\usepackage{geometry}
\geometry{margin=1in}
\usepackage{hyperref}
\hypersetup{colorlinks=true, linkcolor=blue!50!black, citecolor=blue!50!black, urlcolor=blue!50!black}
% \usepackage{microtype}  % disabled: MiKTeX font-expansion environmental issue
\usepackage{graphicx}
\usepackage{booktabs}
\usepackage[numbers,sort&compress]{natbib}
\usepackage{xcolor}

\theoremstyle{plain}
\newtheorem{theorem}{Theorem}[section]
\newtheorem{lemma}[theorem]{Lemma}
\newtheorem{proposition}[theorem]{Proposition}
\newtheorem{corollary}[theorem]{Corollary}
\theoremstyle{definition}
\newtheorem{definition}[theorem]{Definition}
\newtheorem{remark}[theorem]{Remark}
\newtheorem{example}[theorem]{Example}
\newtheorem{observation}[theorem]{Observation}

\newcommand{\disk}{\mathbb{D}}
\newcommand{\diskR}{\disk_{R}}
\newcommand{\cone}{\disk_{R,\alpha}}
\newcommand{\Hilb}{\mathcal{H}}
\newcommand{\Op}{\mathcal{O}}
\newcommand{\Acal}{\mathcal{A}}
\newcommand{\Berez}{\mathcal{B}}
\newcommand{\opnorm}[1]{\lVert #1 \rVert_{\mathrm{op}}}
\newcommand{\supnorm}[1]{\lVert #1 \rVert_{\infty}}
\newcommand{\norm}[1]{\lVert #1 \rVert}
\newcommand{\R}{\mathbb{R}}
\newcommand{\C}{\mathbb{C}}
\newcommand{\N}{\mathbb{N}}
\newcommand{\Z}{\mathbb{Z}}
\newcommand{\Lip}{\mathrm{Lip}}
\newcommand{\MK}{\mathrm{MK}}
\newcommand{\Prob}{\mathcal{P}_{1}}
\newcommand{\Cst}{\mathcal{C}}
\newcommand{\Lebesgue}{\mathcal{L}}
\providecommand{\normone}[1]{\lVert #1 \rVert_{1}}

\title{Spectral truncation of the flat plane:\ state-space convergence
with a squared window, and the disk-with-cone left open}
\author{J.~Loutey\thanks{Independent researcher. \texttt{jloutey@gmail.com}.
This work was developed in an AI-augmented agentic workflow with
Anthropic's Claude. Implementation, internal memos, numerical
verification, and bibliographic collation were performed
collaboratively; all theorems, design choices, and editorial decisions
are the author's responsibility.}}
\date{Preprint; rebuilt 2026-09-29}

\begin{document}
\maketitle

\begin{abstract}
The spectral-truncation convergence results of the companion
Papers~\cite{paper38,paper40} treat compact groups, where a Fej\'er-type
kernel is available and the state spaces are compact. We treat the flat
plane $\R^{2}$ --- non-compact, but with a transitive translation group.
Truncate $L^{2}(\R^{2})$ to the band-limited subspace $\Hilb_{\Lambda}$
(Fourier support $|\xi|\le\sqrt\Lambda$), compress bounded uniformly
continuous multipliers into it, and metrize the truncated state space by
the translation Lipschitz seminorm. With a lift built from the
translates of one window vector $\psi_{\Lambda}\in\Hilb_{\Lambda}$, the
compression/lifted-state argument of Paper~38 (whose scalar
compact-group form is due to Gaudillot-Estrada and
van~Suijlekom~\cite{gaudillot_vs2025}) goes through with the
\emph{positive} kernel $|\psi_{\Lambda}|^{2}$. Four estimates hold with the
single constant $\gamma_{\Lambda}=\int|x|\,|\psi_{\Lambda}(x)|^{2}dx=
\gamma_{1}\Lambda^{-1/2}$ \textbf{[INTERNAL THEOREM]}, and they give a
Gromov--Hausdorff bound $d_{\mathrm{GH}}\le\gamma_{\Lambda}$ between the
Monge--Kantorovich state spaces, taken over the states of finite first
moment \textbf{[INTERNAL THEOREM]} for those state classes;\ whether they
are the canonical ones is open. The window must have a finite first moment.
The sharp ball window does not: its first moment diverges
logarithmically. The smooth-edged windows $(1-|\xi|^{2})_{+}^{m}$ with
$m\ge1$ do; for $m=1$, $\gamma_{1}=32/(5\pi)$. The two-dimensional
Bochner--Riesz map is \emph{not} a usable lift: its kernel is negative
somewhere at every order we tested, and its Lebesgue constant exceeds~$1$
\textbf{[MEASURED]}. On the finite Dirichlet disk-with-cone the paper
reports measurements and withdraws earlier claims. The Markov--Ces\`aro
reconstruction map is not positivity-preserving (a ring counterexample),
and its reconstruction error on the open disk decays on every tested
observable. The ``non-decaying'' rate reported in an earlier version was
an artifact of evaluating a $0/0$ ratio at the boundary point.
Convergence on the disk, and on the cone with apex angle $\alpha\ne1$
(which has no translation group), is left open.
\end{abstract}

% =====================================================================
\section{Introduction}
\label{sec:intro}
% =====================================================================

\subsection{Context}
In the spectral-truncation setting of Connes and
van~Suijlekom~\cite{connes_vs2021} a truncated operator system
$\Op_{\Lambda}=P_{\Lambda}\Acal P_{\Lambda}$ replaces the algebra of
functions, and one asks whether its state space, with a
Monge--Kantorovich metric, converges to the state space of the
continuum. For compact metric groups Gaudillot-Estrada and
van~Suijlekom~\cite{gaudillot_vs2025} prove this with a
compression/lifted-state pair and a single first-moment constant;
Paper~38~\cite{paper38} transports the argument to the spinor window on
$S^{3}=\mathrm{SU}(2)$, and Paper~40~\cite{paper40} to compact semisimple
Lie groups. All of these carriers are compact.

This paper was first written about the two-dimensional disk-with-cone
$\cone$, the spatial factor of the near-horizon cigar of the companion
gravity paper~\cite{paper51}. It claimed a Bochner--Riesz reconstruction
on a ``plane $\R^{2}_{\alpha}$'' and a propinquity backbone for that
paper's cigar-tip entropy. Those claims do not hold, and are withdrawn
(\S\ref{subsec:withdrawn}). What survives is a flat-plane theorem, proved
by a different map, and a set of finite-disk measurements.

\subsection{Results}
\begin{itemize}
\item \textbf{(Flat plane, \S\ref{sec:plane}.)} For the band-limited
truncation of $\R^{2}$ with the translation seminorm, the compression
$S$ and the window lift $\upsilon$ satisfy the four estimates of
Paper~38's unconditional proof with the positive kernel
$|\psi_{\Lambda}|^{2}$ and $\gamma_{\Lambda}=\gamma_{1}\Lambda^{-1/2}$
(Lemma~\ref{lem:four_facts}, \textbf{[INTERNAL THEOREM]}). Consequently
(Theorem~\ref{thm:plane_gh}, \textbf{[INTERNAL THEOREM]}) the Gromov--Hausdorff
distance between the finite-first-moment state classes is at most
$\gamma_{\Lambda}$. Unlike the Bochner--Riesz route, no Lebesgue constant
enters, because the kernel is positive.
\item \textbf{(Window choice, \S\ref{subsec:windows}.)} The windows
$\widehat\psi_{1}\propto(1-|\xi|^{2})_{+}^{m}$ give
$|\psi_{1}|^{2}=c_{m}J_{m+1}(r)^{2}/r^{2m+2}$, the squared modulus of the
order-$m$ Bochner--Riesz kernel. Its first moment is finite exactly when
$m\ge1$; $\gamma_{1}=32/(5\pi)\approx2.04$ at $m=1$
\textbf{[INTERNAL THEOREM + MEASURED]}. The sharp window ($m=0$) has a
logarithmically divergent first moment \textbf{[INTERNAL THEOREM +
MEASURED]}.
\item \textbf{(Why not Bochner--Riesz, \S\ref{subsec:not_br}.)} The
two-dimensional Bochner--Riesz kernel is negative somewhere at
$s=0.75,1,2,4$ (min/max $=-0.071,-0.059,-0.029,-0.009$), and its
$\Lambda$-independent Lebesgue constant exceeds $1$ at every finite order
\textbf{[MEASURED]}.
\item \textbf{(Finite disk, \S\ref{sec:disk}.)} On the Dirichlet disk the
Markov--Ces\`aro map is unital. It is \emph{not}
positivity-preserving: a non-negative ring observable is sent to a function
with negative values, at Ces\`aro orders $2$ and $4$. Its reconstruction
error on the open disk decays on every tested observable. The operator
system has propagation number two at the four truncations computed. All
of these are \textbf{[MEASURED]}; no disk convergence theorem is claimed.
\end{itemize}

\subsection{Honest scope}\label{subsec:withdrawn}
Withdrawn in the 2026-09-29 rebuild, each for the reason stated:
\begin{enumerate}
\item \emph{Positivity of the two-dimensional Bochner--Riesz means.}
False: the kernel oscillates (\S\ref{subsec:not_br}). The earlier test
checked one positive Gaussian, which does not test positivity
preservation.
\item \emph{A ``plane $\R^{2}_{\alpha}$'' reconstruction and a backbone for
the cigar-tip entropy of Paper~51.} The earlier plane map was defined
through the order-$0$ Hankel transform, i.e.\ on the radial sector of
the \emph{flat} plane; the apex angle $\alpha$ never entered. The cone
with $\alpha\ne1$ has no translation group, so nothing in this paper
reaches it (\S\ref{sec:cone}). Paper~51 cites this paper only for the
withdrawn material;\ its Theorem on the cigar-tip entropy does not depend on
it.
\item \emph{Plane van~Suijlekom state-space convergence via a
Bochner--Riesz tunnel.} Not established: its ``height'' leg carried a
Lebesgue constant $>1$ and the constituents did not match the framework
(Remark~\ref{rem:height_constant}). It is replaced by
Theorem~\ref{thm:plane_gh}, proved by a different map.
\item \emph{The finite-disk ``non-decaying interior rate''
$\Lambda^{+0.07}$ and the ``Dirichlet-boundary obstruction''.} The
plateau near $0.3$ came from evaluating the Markov ratio at the exact
grid point $\rho=R$, where numerator and denominator are both
floating-point rounding noise. On the open disk the error decays
(\S\ref{sec:disk}). The boundary-collar error ``$\approx0.26$ at $R=12$''
is not reproduced: a direct recomputation gives $\sim10^{-11}$.
\item \emph{Ces\`aro positivity on the finite disk as a general
property.} It held on the single tested observable $(\rho/R)^{2}$ and
fails on a non-negative ring (\S\ref{sec:disk}).
\end{enumerate}
Not claimed: any statement about Latr\'emoli\`ere's
propinquity~\cite{latremoliere2015,latremoliere2016,latremoliere2025},
any convergence on the finite disk or the cone, and novelty of the
proof pattern (Remark~\ref{rem:prior_art}).

% =====================================================================
\section{The flat plane}
\label{sec:plane}
% =====================================================================

\subsection{Setup}
Let $\Hilb=L^{2}(\R^{2})$, with unitary translations
$(U_{a}\phi)(x)=\phi(x-a)$, $a\in\R^{2}$. For $\Lambda>0$ let $P_{\Lambda}$
be the Fourier projection onto $\Hilb_{\Lambda}=\{\phi:\operatorname{supp}
\widehat\phi\subset\{|\xi|\le\sqrt\Lambda\}\}$. Translations are Fourier
multipliers, so $[U_{a},P_{\Lambda}]=0$.

Let $\Acal=C_{\mathrm{ub}}(\R^{2})$, the bounded uniformly continuous
functions (a unital C$^{*}$-algebra on which $a\mapsto\tau_{a}f:=f(\cdot-a)$
is norm-continuous), acting by multiplication operators $M_{f}$. Then
$U_{a}M_{f}U_{a}^{*}=M_{\tau_{a}f}$. The truncated operator system is the
norm closure
\[
\Op_{\Lambda}=\overline{\{S(f):=P_{\Lambda}M_{f}P_{\Lambda}\,:\,f\in\Acal\}}
\subset B(\Hilb_{\Lambda}),
\]
unital since $S(1)=P_{\Lambda}$ is the identity of $B(\Hilb_{\Lambda})$.
$S:\Acal\to\Op_{\Lambda}$ is unital and completely positive.

\begin{definition}[translation seminorms]\label{def:translation}
For $f\in\Acal$ and $T\in\Op_{\Lambda}$,
\[
L(f)=\sup_{a\ne0}\frac{\supnorm{\tau_{a}f-f}}{|a|},\qquad
L_{\Lambda}(T)=\sup_{a\ne0}\frac{\opnorm{U_{a}TU_{a}^{*}-T}}{|a|}.
\]
\end{definition}
Since translations act transitively by isometries of the Euclidean
metric, $L(f)=\Lip(f)$, and $L(f)=\supnorm{\nabla f}$ for $f\in C^{1}$.
On $\Hilb_{\Lambda}$, $\opnorm{U_{a}-1}\le|a|\sqrt\Lambda$ (the symbol
$e^{-ia\cdot\xi}$ on the ball), so $L_{\Lambda}(T)\le2\sqrt\Lambda
\opnorm{T}<\infty$ for every $T$.

\subsection{The window lift}
\begin{definition}[window and lift]\label{def:window}
Let $\varphi\in L^{2}(\R^{2})$ be supported in the closed unit ball, with
$\norm{\varphi}_{2}=1$, and let $\psi_{1}=\mathcal F^{-1}\varphi$
(unitary Fourier transform). Assume
\begin{equation}\label{eq:gamma1}
\gamma_{1}:=\int_{\R^{2}}|x|\,|\psi_{1}(x)|^{2}\,dx<\infty .
\end{equation}
Set $\psi_{\Lambda}(x)=\Lambda^{1/2}\psi_{1}(\Lambda^{1/2}x)$, so that
$\norm{\psi_{\Lambda}}_{2}=1$ and $\widehat{\psi_{\Lambda}}$ is supported
in $|\xi|\le\sqrt\Lambda$; hence $U_{x}\psi_{\Lambda}\in\Hilb_{\Lambda}$
for every $x$. Write $k_{\Lambda}=|\psi_{\Lambda}|^{2}$, a probability
density, and define
\[
\upsilon(T)(x)=\langle U_{x}\psi_{\Lambda},\,T\,U_{x}\psi_{\Lambda}\rangle,
\qquad
\gamma_{\Lambda}=\int|x|\,k_{\Lambda}(x)\,dx .
\]
\end{definition}

\begin{lemma}[scaling; \textbf{[INTERNAL THEOREM]}]\label{lem:scaling}
$\gamma_{\Lambda}=\gamma_{1}\Lambda^{-1/2}$.
\end{lemma}
\begin{proof}
Substitute $y=\Lambda^{1/2}x$ in
$\int|x|\,\Lambda|\psi_{1}(\Lambda^{1/2}x)|^{2}dx$.
\end{proof}

\begin{lemma}[four estimates; \textbf{[INTERNAL THEOREM]}]\label{lem:four_facts}
For $f\in\Acal$ and $T\in\Op_{\Lambda}$:
\begin{enumerate}\setlength{\itemsep}{1pt}
\item[(i)] $L_{\Lambda}(S(f))\le L(f)$;
\item[(ii)] $\upsilon(T)\in\Acal$, $\upsilon$ is unital and positive, and
$L(\upsilon(T))\le L_{\Lambda}(T)$;
\item[(iii)] $\upsilon(S(f))(x)=\int f(x+z)\,k_{\Lambda}(z)\,dz$, hence
$\supnorm{\upsilon(S(f))-f}\le\gamma_{\Lambda}L(f)$;
\item[(iv)] $S(\upsilon(T))=\Phi(T):=\int k_{\Lambda}(z)\,U_{-z}TU_{-z}^{*}\,dz$
(weak integral), hence $\opnorm{S(\upsilon(T))-T}\le\gamma_{\Lambda}
L_{\Lambda}(T)$.
\end{enumerate}
\end{lemma}

\begin{proof}
(i) $U_{a}S(f)U_{a}^{*}-S(f)=P_{\Lambda}M_{\tau_{a}f-f}P_{\Lambda}$, and
compression does not increase norms.

(ii) $|\upsilon(T)|\le\opnorm{T}$, and $\upsilon(1)=1$, $\upsilon(T)\ge0$
for $T\ge0$, from the vector-state form. Since $U_{x}^{*}=U_{-x}$,
\[
\upsilon(T)(x)-\upsilon(T)(x')=\bigl\langle\psi_{\Lambda},\,
U_{-x'}\bigl(U_{x'-x}TU_{x'-x}^{*}-T\bigr)U_{-x'}^{*}\,\psi_{\Lambda}
\bigr\rangle,
\]
whose modulus is at most $L_{\Lambda}(T)|x-x'|$. So $\upsilon(T)$ is
bounded and Lipschitz, hence in $\Acal$, with $L(\upsilon(T))\le
L_{\Lambda}(T)$.

(iii) $U_{x}\psi_{\Lambda}\in\Hilb_{\Lambda}$ absorbs the projections:
$\upsilon(S(f))(x)=\int f(y)|\psi_{\Lambda}(y-x)|^{2}dy=\int f(x+z)
k_{\Lambda}(z)dz$. Then $|\upsilon(S(f))(x)-f(x)|\le\int|f(x+z)-f(x)|
k_{\Lambda}(z)dz\le L(f)\gamma_{\Lambda}$.

(iv) For $T=S(f)$, (iii) gives $M_{\upsilon(T)}=\int k_{\Lambda}(z)\,
M_{f(\cdot+z)}\,dz$ weakly, and $M_{f(\cdot+z)}=U_{-z}M_{f}U_{-z}^{*}$;
compressing, and using $[U_{-z},P_{\Lambda}]=0$, gives
$S(\upsilon(T))=\Phi(T)$. Both sides are norm-contractive and linear in
$T$, so the identity extends to the closure $\Op_{\Lambda}$. Finally
$\opnorm{\Phi(T)-T}\le\int k_{\Lambda}(z)\opnorm{U_{-z}TU_{-z}^{*}-T}dz\le
L_{\Lambda}(T)\int|z|k_{\Lambda}(z)dz$.
\end{proof}

The positivity of $k_{\Lambda}$ is used twice: in (iii) and (iv) the
average of a quantity bounded by $L\cdot|z|$ is bounded by $L\gamma_{\Lambda}$.
With a signed kernel of $L^{1}$ norm $\Lebesgue>1$ the same steps give no
such bound (a signed kernel is not a state-induced average; compare
Remark~\ref{rem:height_constant}).

\begin{corollary}[kernel; \textbf{[INTERNAL THEOREM]}]\label{cor:kernel}
$L_{\Lambda}(T)=0$ if and only if $T\in\C1$.
\end{corollary}
\begin{proof}
If $L_{\Lambda}(T)=0$ then by (iv) $T=S(\upsilon(T))$, and by (ii)
$\upsilon(T)$ is constant.
\end{proof}

\subsection{State classes and the convergence statement}
The Monge--Kantorovich (MK) metric on all states of $\Acal$ is not
finite. For example $f_{n}(x)=\min(|x|,n)$ has $L(f_{n})=1$ and separates
$\delta_{0}$ from $\delta_{x}$ by $\min(|x|,n)$. Since $\Acal$ also has
non-normal states at infinity, we restrict to states of finite first
moment on both sides:
\[
\Prob=\{\mu:\ \text{Borel probability measure on }\R^{2},\
\textstyle\int|x|\,d\mu<\infty\},\qquad
\Cst_{\Lambda}=\{\omega\in S(\Op_{\Lambda}):\ S^{*}\omega\in\Prob\},
\]
where $S^{*}\omega=\omega\circ S$ and ``$\in\Prob$'' means that it is
integration against such a measure. On $\Prob$ the MK metric of $L$ is
the Wasserstein distance $W_{1}$ (Kantorovich--Rubinstein
duality~\cite{villani2009}; bounded Lipschitz test functions suffice,
by truncation, when first moments are finite). On $\Cst_{\Lambda}$ put
$\MK_{\Lambda}(\omega,\omega')=\sup\{|\omega(T)-\omega'(T)|:T=T^{*}\in
\Op_{\Lambda},\ L_{\Lambda}(T)\le1\}$.

\begin{theorem}[flat-plane state-space convergence]\label{thm:plane_gh}
Under Definition~\ref{def:window},
\[
d_{\mathrm{GH}}\Bigl(\bigl(\Cst_{\Lambda},\MK_{\Lambda}\bigr),\
\bigl(\Prob,W_{1}\bigr)\Bigr)\ \le\ \gamma_{\Lambda}
=\gamma_{1}\Lambda^{-1/2}\ \xrightarrow[\Lambda\to\infty]{}\ 0 .
\]
\textbf{Tiers.} The four estimates, the scaling and the kernel
statement are \textbf{[INTERNAL THEOREM]}
(Lemmas~\ref{lem:scaling}--\ref{lem:four_facts},
Corollary~\ref{cor:kernel}). The Gromov--Hausdorff inequality is
\textbf{[INTERNAL THEOREM]} for the two metric spaces named:\ the proof below
closes on the classes $\Cst_{\Lambda}$, $\Prob$. We have not shown that these
classes are the canonical ones (an open question, not a gap in the proof), and we have not connected them to Latr\'emoli\`ere's
pointed/proper framework~\cite{latremoliere2025}.
\end{theorem}

\begin{proof}
\emph{Well-definedness.} For $\mu\in\Prob$ put $\upsilon^{*}\mu(T)=
\int\upsilon(T)\,d\mu$, a state of $\Op_{\Lambda}$ by
Lemma~\ref{lem:four_facts}(ii). By (iii), $S^{*}\upsilon^{*}\mu=\mu*k_{\Lambda}$
(for the radial windows used $k_{\Lambda}$ is even, so the reflection
convention is immaterial), whose first moment is at
most $\int|x|d\mu+\gamma_{\Lambda}<\infty$; so $\upsilon^{*}\mu\in
\Cst_{\Lambda}$. By definition $S^{*}$ maps $\Cst_{\Lambda}$ into $\Prob$.
$\MK_{\Lambda}$ is finite on $\Cst_{\Lambda}$: for $L_{\Lambda}(T)\le1$,
(iv) and (ii) give $|\omega(T)-\omega'(T)|\le W_{1}(S^{*}\omega,
S^{*}\omega')+2\gamma_{\Lambda}$. It separates states, since every
self-adjoint $T$ has $L_{\Lambda}(T)<\infty$.

\emph{Lipschitz pullbacks.} By (i), $W_{1}(S^{*}\omega,S^{*}\omega')\le
\MK_{\Lambda}(\omega,\omega')$; by (ii), $\MK_{\Lambda}(\upsilon^{*}\mu,
\upsilon^{*}\nu)\le W_{1}(\mu,\nu)$.

\emph{Almost inverses.} By (iii), $W_{1}(S^{*}\upsilon^{*}\mu,\mu)=
\sup_{L(f)\le1}|\int(\upsilon S f-f)d\mu|\le\gamma_{\Lambda}$; by (iv),
$\MK_{\Lambda}(\upsilon^{*}S^{*}\omega,\omega)=\sup_{L_{\Lambda}(T)\le1}
|\omega(S\upsilon T-T)|\le\gamma_{\Lambda}$.

\emph{Correspondence.} Let $\mathcal R=\{(\mu,\upsilon^{*}\mu)\}\cup
\{(S^{*}\omega,\omega)\}$. For two pairs of the first kind,
$\MK_{\Lambda}(\upsilon^{*}\mu,\upsilon^{*}\nu)\le W_{1}(\mu,\nu)\le
W_{1}(S^{*}\upsilon^{*}\mu,S^{*}\upsilon^{*}\nu)+2\gamma_{\Lambda}\le
\MK_{\Lambda}(\upsilon^{*}\mu,\upsilon^{*}\nu)+2\gamma_{\Lambda}$. For
a mixed pair, $W_{1}(\mu,S^{*}\omega)\le\MK_{\Lambda}(\upsilon^{*}\mu,
\omega)+\gamma_{\Lambda}$ and $\MK_{\Lambda}(\upsilon^{*}\mu,\omega)\le
W_{1}(\mu,S^{*}\omega)+\gamma_{\Lambda}$. The case of two pairs of the
second kind is symmetric. The distortion is therefore at most
$2\gamma_{\Lambda}$, and $d_{\mathrm{GH}}\le\gamma_{\Lambda}$. No
compactness is used; both spaces are unbounded, and the bound is on the
full (not pointed) Gromov--Hausdorff distance.
\end{proof}

\subsection{Windows: smooth edge versus sharp edge}\label{subsec:windows}
\begin{proposition}[the $(1-|\xi|^{2})^{m}$ windows;\
\textbf{[INTERNAL THEOREM + MEASURED]}]\label{prop:windows}
Let $\varphi_{m}\propto(1-|\xi|^{2})_{+}^{m}$, $m\in\{0,1,2,\dots\}$. By Sonine's
integral~\cite{watson1944},
\[
|\psi_{1}(x)|^{2}=c_{m}\,\frac{J_{m+1}(r)^{2}}{r^{2m+2}},\qquad
r=|x|,\qquad c_{m}=\frac{4^{m}(m!)^{2}(2m+1)}{\pi},
\]
which is the squared modulus of the two-dimensional Bochner--Riesz kernel
of order $m$. It decays like $r^{-2m-3}$, so \eqref{eq:gamma1} holds if and only if
$m\ge1$. For $m=1$, $\gamma_{1}=32/(5\pi)\approx2.037$; numerically
$\gamma_{1}\approx2.299$ ($m=2$) and $2.593$ ($m=3$).
\end{proposition}
The normalisation $c_{m}$ and the $m=1$ moment follow from the
Weber--Schafheitlin-type integral $\int_{0}^{\infty}J_{\nu}(t)^{2}
t^{-\lambda}dt$. All three moments are reproduced by quadrature in
\texttt{tests/test\_paper53\_flat\_plane\_squared\_window.py}. The $\pi$
in $\gamma_{1}$ is the normalisation of the continuum Fourier--Bessel
transform, a continuum-side constant. It is not a skeleton quantity, and
no Paper~34 catalogue placement is claimed.

\begin{remark}[the sharp window diverges]\label{rem:sharp}
\textbf{[INTERNAL THEOREM + MEASURED]} For $m=0$ (indicator of the ball),
$k_{1}=J_{1}(r)^{2}/(\pi r^{2})$. Its partial first moment is
$\int_{|x|\le\rho}|x|k_{1}=2\int_{0}^{\rho}J_{1}(t)^{2}dt$, which grows
like $(2/\pi)\log\rho$ by the Bessel asymptotics; the test measures the
slope. So in two dimensions the sharp spectral window gives no finite
rate, and a smooth-edged window is needed. On the compact group
$\mathrm{SU}(2)$ the sharp (Fej\'er) window does give a finite rate,
$(4/\pi)\log n/n$ in Paper~38's normalisation. A natural reading is that
the same slowly decaying tail is cut off there at the diameter. That
reading is \textbf{[ARGUED]}; we have not checked it here.
\end{remark}

\subsection{Why the Bochner--Riesz map is not the lift}\label{subsec:not_br}
The earlier version of this paper used the two-dimensional Bochner--Riesz
mean
\begin{equation}\label{eq:plane_berezin}
\Berez_{\Lambda}^{s}(f)=\mathcal F^{-1}\bigl[(1-|\xi|^{2}/\Lambda)_{+}^{s}
\,\widehat f\,\bigr]=K_{s,\Lambda}*f,\qquad
K_{s,1}(r)\propto\frac{J_{s+1}(r)}{r^{s+1}},
\end{equation}
as its reconstruction map. Two measured facts rule it out as a lift.

\emph{It does not preserve positivity.} \textbf{[MEASURED]} $K_{s,1}$
oscillates: $\min K/\max K=-0.071,-0.059,-0.029,-0.009$ at
$s=0.75,1,2,4$ (\texttt{tests/test\_paper53\_plane\_bochner\_riesz.py}).
A non-negative $f$ concentrated where $K<0$ is therefore mapped to a
function with negative values. The positivity claimed in the earlier
version (``classical, $s\ge\tfrac12$'') is false. Bochner--Riesz summability
above the critical index~\cite{stein_weiss} concerns convergence, not
positivity.

\begin{remark}[the Lebesgue constant exceeds $1$ --- corrected 2026-09-07]
\label{rem:height_constant}
\textbf{[MEASURED]} $\Berez_{\Lambda}^{s}$ is a convolution, so
$\nabla(\Berez f)=\Berez(\nabla f)$, and by Young's inequality the
sharp gradient gain is the Lebesgue constant
$\Lebesgue_{s}:=\normone{K_{s}}/\int K_{s}$. This equals $1$ only for a
non-negative kernel. Integrating over Bessel half-waves:
\[
\begin{array}{l|cccccccc}
s & 0.6 & 0.75 & 1 & 1.5 & 2 & 3 & 5 & 8\\\hline
\Lebesgue_{s} & 5.72 & 3.23 & 2.01 & 1.42 & 1.23 & 1.09 & 1.02 & 1.004
\end{array}
\]
$\Lebesgue_{s}$ is $\Lambda$-independent by scaling and tends to $1$
only as $s\to\infty$ (\texttt{tests/test\_paper53\_height\_constant.py}).
An earlier version asserted that $\Berez_{\Lambda}^{s}$ is
gradient-non-expansive, ``ratio $\le1$ at every $\Lambda$''. That is false
over the unit-Lipschitz ball. The $\Lambda$-dependence it reported was a
property of the Gaussian test functions. This is not Paper~38's
withdrawn Lemma~L5 transported:\ L5's height was a reconstruction-defect
quantity of a different construction.

\textbf{Consequence for the construction.} In the argument of
Lemma~\ref{lem:four_facts}, a lift of the form
$T\mapsto K*(\text{something})$ with a signed kernel of $L^{1}$ norm
$\Lebesgue_{s}>1$ is not a unital positive map. The ``height'' used
in the earlier assembly (a Lipschitz gain) was also not Latr\'emoli\`ere's
height, which is a state-space Hausdorff distance required to
vanish~\cite{latremoliere_metric_st,
latremoliere2025}. Both problems disappear with the squared window. Its
kernel $|\psi_{\Lambda}|^{2}$, the squared modulus of a Bochner--Riesz
kernel when $\widehat\psi\propto(1-|\xi|^{2})^{m}_{+}$
(Proposition~\ref{prop:windows}), is positive, so the constant is exactly
$1$ and the only rate is the first moment $\gamma_{\Lambda}$.
\end{remark}

\begin{remark}[the Bochner--Riesz rate on smooth observables]
\label{rem:br_rate}
\textbf{[MEASURED]} On individual smooth observables the Bochner--Riesz
mean does reconstruct. For $f=e^{-\rho^{2}/2}$ at $s=2$,
$\supnorm{\Berez_{\Lambda}^{s}f-f}/\supnorm{\nabla f}$ decays roughly like
$\Lambda^{-0.9}$ ($\Lambda=5\to160$), and like $\Lambda^{-0.6}$ for a
narrower Gaussian (\texttt{tests/test\_paper53\_plane\_bochner\_riesz.py}).
These are observable-dependent rates. The uniform rate over the
Lipschitz ball, which is what the Gromov--Hausdorff bound needs, is the
first-moment rate $\gamma_{\Lambda}\propto\Lambda^{-1/2}$ of
Theorem~\ref{thm:plane_gh}.
\end{remark}

\begin{remark}[prior art]\label{rem:prior_art}
The compression/lifted-state pattern, with a single first-moment
constant, is Gaudillot-Estrada and van~Suijlekom's for compact metric
groups~\cite{gaudillot_vs2025}; Paper~38~\cite{paper38} uses it on the
$S^{3}$ spinor window. The present adaptation replaces the group by the
non-compact translation group of $\R^{2}$, the Fej\'er kernel by a
squared band-limited window, and the compact state spaces by
finite-first-moment classes. We did not find a treatment of flat
$\R^{n}$ spectral truncations of this kind in the literature we checked.
This is a statement about our search, not a novelty claim.
\end{remark}

% =====================================================================
\section{The finite disk-with-cone: measurements}
\label{sec:disk}
% =====================================================================

\subsection{Setup}
Let $\diskR$ be the flat disk of radius $R$ with the Dirichlet Laplacian,
eigenbasis
\begin{equation}
\psi_{k,j}(\rho,\phi)=N_{k,j}\,J_{|k|}\!\Big(\frac{a_{|k|,j}\rho}{R}\Big)
\frac{e^{ik\phi}}{\sqrt{2\pi}},\qquad
\lambda_{k,j}=\Big(\frac{a_{|k|,j}}{R}\Big)^{2},
\label{eq:eigenbasis}
\end{equation}
with $a_{m,j}$ the $j$-th positive zero of $J_{m}$ and
$N_{k,j}=\sqrt{2}/(R\,|J_{|k|+1}(a_{|k|,j})|)$. On the disk-with-cone
$\cone$ (azimuthal period $2\pi\alpha$) the half-integer fermionic
azimuthal momentum is $m_{\mathrm{eff}}=(k+\tfrac12)/\alpha$. Let
$P_{\Lambda}$ project onto $\lambda\le\Lambda$ and
$\Op_{\Lambda}=P_{\Lambda}C(\cone)P_{\Lambda}$.

\begin{observation}[propagation number two; \textbf{[MEASURED]}]
\label{prop:prop2}
In the truncated radial$\times$azimuthal mode basis at $\alpha=1$ (fermionic
$m=|k+\tfrac12|$), built from compressed monomials $z^{a}\bar z^{b}$, the
squared operator system $\Op^{2}$ spans the full envelope $M_{N}(\C)$ at
$(K,J)\in\{(2,2),(2,3),(3,2),(3,3)\}$, $N\in\{8,12,12,18\}$
(\texttt{tests/test\_paper53\_disk\_obstruction.py}). This matches the
Connes--van~Suijlekom Toeplitz value~\cite{connes_vs2021}. It is not proved
for all cutoffs.
\end{observation}

\subsection{The Markov--Ces\`aro map}
\begin{definition}[Markov--Ces\`aro reconstruction]\label{def:berezin}
For order $s$ and cutoff $\Lambda$, with $w_{a}=(1-\lambda_{a}/\Lambda)_{+}^{s}$,
\begin{equation}
g_{\Lambda}(x)=
\frac{\sum_{a}w_{a}\langle\psi_{a},f\rangle\psi_{a}(x)}
{\sum_{a}w_{a}\langle\psi_{a},1\rangle\psi_{a}(x)},
\qquad
\Berez_{\Lambda}(f)=P_{\Lambda}M_{g_{\Lambda}}P_{\Lambda}.
\label{eq:berezin}
\end{equation}
\end{definition}

\begin{proposition}[what holds and what does not; \textbf{[MEASURED]}
except (a)]\label{prop:disk}
In the radial ($k=0$) sector, $R=1$, $\Lambda=100$ to $1600$ (and
$6400$ where stated):
\begin{enumerate}
\item[(a)] \emph{Unitality} (exact): $f\equiv1\Rightarrow g_{\Lambda}\equiv1$.
\item[(b)] \emph{Positivity fails in general.} For a non-negative ring
$f(\rho)=e^{-((\rho-\rho_{0})/0.02)^{2}}$, centred where the Ces\`aro
kernel is most negative, $\min g_{\Lambda}=-0.016,-0.025,-0.044$ at
$\Lambda=100,400,1600$ ($s=2$). At $s=4$ it is $-0.006$ and $-0.011$ at
$\Lambda=400,1600$. The denominator is positive at those points. The
interior Ces\`aro kernel has $\min/\max\approx-0.03$ ($s=2$) and
$-0.009$ ($s=4$), matching the plane Bochner--Riesz values of
\S\ref{subsec:not_br}. On the single observable $(\rho/R)^{2}$, $\min
g_{\Lambda}>0$ at $s=2$; this is where the earlier general positivity
claim came from. The truncated heat weight $e^{-\lambda/\Lambda}$ already
fails on $(\rho/R)^{2}$ ($\min g\approx-0.13$ at $\Lambda=200$).
\item[(c)] \emph{The reconstruction error on the open disk decays.} For
$f=(\rho/R)^{2}$, which is non-zero at the boundary,
$\sup_{\rho<R}|g_{\Lambda}-f|/\supnorm{\nabla f}=0.237,0.179,0.132,
0.097,0.070$ at $\Lambda=100,200,400,800,1600$ (about $\Lambda^{-0.44}$). The
maximum sits at the boundary. For an interior-supported $C^{1}$ bump the
error falls roughly like $\Lambda^{-0.9}$ ($0.163\to0.0032$ over
$\Lambda=100\to6400$).
\end{enumerate}
\end{proposition}

\begin{remark}[the withdrawn ``Dirichlet-boundary obstruction'']
\label{rem:interior_correction}
Earlier versions reported that the disk error does not decay
($\sim\Lambda^{+0.07}$, plateau $\approx0.3$) and read this as an
obstruction intrinsic to the Dirichlet boundary. The plateau came from
the single grid point $\rho=R$. There every $J_{0}(a_{j})$ is rounding
noise of order $10^{-16}$, so the ratio~\eqref{eq:berezin} returns an
arbitrary number. Dropping that one point, or evaluating at $R-10^{-9}$,
gives the decaying values of Proposition~\ref{prop:disk}(c) (the test
now pins both). The companion claim that the boundary-collar error for
$f=e^{-\rho^{2}/2}$ on $R=12$ stays near $0.26$ is not reproduced: a
direct recomputation gives $\sim10^{-11}$, in line with the sprint
driver's own docstring. What the boundary does cause is a slower rate for
observables that are non-zero there. The disk measurements therefore do
\emph{not} show an obstruction. They show that the obvious reconstruction
map fails positivity, and that no proof route is available.
\end{remark}

\begin{remark}[why no disk theorem]
The argument of \S\ref{sec:plane} needs a transitive group of isometries
commuting with the truncation, and a positive lift. The disk has neither:
its isometries are the rotations about the centre, and the
Markov--Ces\`aro map is not positive. Whether the disk truncations
converge in any state-space sense is open (\S\ref{sec:open}, Q2).
\end{remark}

% =====================================================================
\section{The cone $\alpha\ne1$}
\label{sec:cone}
% =====================================================================
The cone $\R^{2}_{\alpha}$ (azimuthal period $2\pi\alpha$) is flat away
from its apex, but for $\alpha\ne1$ it is not homogeneous. Its only
isometries are rotations about the apex, so there is no translation
seminorm and no family of window translates to average over. The earlier
``plane $\R^{2}_{\alpha}$'' construction never reached the cone: it used
the order-$0$ Hankel transform, the radial sector of the flat plane, and
$\alpha$ did not appear. Convergence of spectral truncations of the cone,
and in particular anything bearing on the apex-local cigar-tip entropy of
Paper~51, is open (Q1).

% =====================================================================
\section{Open questions}
\label{sec:open}
% =====================================================================

\paragraph{Q1 (the cone).} Is there a state-space convergence statement
for band-limited (Hankel-order $|k+\tfrac12|/\alpha$) truncations of
the cone $\R^{2}_{\alpha}$? A lift would have to be built without
translations. One option is window vectors localised along geodesics,
with a separate estimate near the apex.

\paragraph{Q2 (the finite disk).} Does a positive, unital reconstruction
map with a decaying Lipschitz-uniform rate exist for the Dirichlet disk?
Such a map would, for instance, be built from localised vectors in the
truncated space rather than from Ces\`aro weights. And is the propagation
number two for all cutoffs (Observation~\ref{prop:prop2})?

\paragraph{Q3 (optimal window and sharp constant).} Among admissible
windows the rate is $\gamma_{1}\Lambda^{-1/2}$. Which window minimises
$\gamma_{1}$, and is $\Lambda^{-1/2}$ the sharp order of the
Gromov--Hausdorff distance, or only of this bound?

\paragraph{Q4 (Latr\'emoli\`ere's propinquity).} Theorem~\ref{thm:plane_gh}
is a Gromov--Hausdorff statement on chosen state classes. Whether the
same pair $(S,\upsilon)$ yields convergence in the pointed/proper
propinquity~\cite{latremoliere2025} is open.

\paragraph{Q5 ($\R^{n}$).} The argument is dimension-independent in form,
with $\gamma_{\Lambda}=\gamma_{1}\Lambda^{-1/2}$ in every dimension.
Which windows have a finite first moment in dimension $n$ was not checked.

% =====================================================================
\section{Conclusion}
% =====================================================================
On the flat plane, spectral truncation converges in the state-space
Gromov--Hausdorff sense at rate $\gamma_{1}\Lambda^{-1/2}$, provided the
lift uses a window whose squared modulus has a finite first moment. The
four estimates behind this are proved; the final Gromov--Hausdorff step
is argued on finite-first-moment state classes that we chose. The
Bochner--Riesz map is not a usable lift, because its kernel changes sign.
Replacing it by the squared modulus of such a kernel removes the
Lebesgue constant. On the finite disk-with-cone the earlier
obstruction was a numerical artifact, and the obvious reconstruction map
is not positive. The disk and the cone with $\alpha\ne1$ are open.

% =====================================================================
\section*{Acknowledgments}
% =====================================================================
The finite-disk measurements and the Lebesgue-constant table were made in
the GeoVac repository and are frozen in the \texttt{tests/test\_paper53\_*}
files cited above. The 2026-09-29 rebuild, including the retractions in
\S\ref{subsec:withdrawn}, is recorded in the CHANGELOG.

% =====================================================================
\begin{thebibliography}{99}
% =====================================================================

\bibitem{gaudillot_vs2025}
Y.~Gaudillot-Estrada and W.~D.~van~Suijlekom, ``Convergence of spectral
truncations for compact metric groups,'' \textit{Int.~Math.~Res.~Not.\
IMRN} \textbf{2025}, no.~13 (2025), rnaf197;\ arXiv:2310.14733.

\bibitem{connes_vs2021}
A.~Connes and W.~D.~van~Suijlekom, ``Spectral truncations in
noncommutative geometry and operator systems,''
\textit{Comm.~Math.~Phys.}\ \textbf{383} (2021), 2021--2067.

\bibitem{latremoliere2015}
F.~Latr\'emoli\`ere, ``The dual Gromov--Hausdorff propinquity,''
\textit{J.~Math.~Pures Appl.}\ \textbf{103} (2015), 303--351.

\bibitem{latremoliere2016}
F.~Latr\'emoli\`ere, ``Quantum metric spaces and the Gromov--Hausdorff
propinquity,'' \textit{Contemp.~Math.}\ \textbf{676} (2016), 47--133.

\bibitem{latremoliere_metric_st}
F.~Latr\'emoli\`ere, ``The Gromov--Hausdorff propinquity for metric
spectral triples,'' \textit{Adv.~Math.}\ \textbf{404} (2022), Paper
No.~108393, 56\,pp.;\ preprint arXiv:1811.10843.

\bibitem{latremoliere2025}
F.~Latr\'emoli\`ere, ``The quantum Gromov--Hausdorff hypertopology on
the class of pointed proper quantum metric spaces,'' arXiv:2512.03573
(2025).

\bibitem{villani2009}
C.~Villani, \textit{Optimal Transport: Old and New}, Grundlehren der
mathematischen Wissenschaften \textbf{338}, Springer, 2009
(Kantorovich--Rubinstein duality).

\bibitem{watson1944}
G.~N.~Watson, \textit{A Treatise on the Theory of Bessel Functions},
2nd ed., Cambridge Univ.\ Press, 1944 (Sonine's finite integral;
Weber--Schafheitlin integrals).

\bibitem{stein_weiss}
E.~M.~Stein and G.~Weiss, \textit{Introduction to Fourier Analysis on
Euclidean Spaces}, Princeton Univ.\ Press, 1971 (Bochner--Riesz means).

\bibitem{paper38}
GeoVac Project. Paper~38 --- State-space Gromov--Hausdorff convergence
of truncated Camporesi--Higuchi spectral triples on $S^3$.
\texttt{papers/group1\_operator\_algebras/paper\_38\_su2\_propinquity\_convergence.tex}.

\bibitem{paper40}
GeoVac Project. Paper~40 --- State-space Gromov--Hausdorff convergence of
spectral truncations of compact semisimple Lie groups with bi-invariant metric.
\texttt{papers/group1\_operator\_algebras/paper\_40\_unified\_propinquity\_convergence.tex}.

\bibitem{paper51}
GeoVac Project. Paper~51 --- Gravity from the GeoVac spectral action:
continuum results and the discrete-substrate program.
\texttt{papers/group5\_qed\_gauge/paper\_51\_gravity\_arc.tex}.

\end{thebibliography}

\end{document}
